\section{Formulación matemática}\label{app:formulacion}
Este anexo desarrolla la formulación de la Sección~\ref{sec:formulacion}: el problema de optimización, el
modelo de escenarios y de valorización que define sus coeficientes, las medidas que se reportan y los
supuestos y decisiones de modelación en que se apoya. Los montos están en S/ millones, las tasas en
decimales y el tiempo en años ACT/365.

\subsection{Notación}
Sea $t_0$ la fecha de decisión y $t_H$ el horizonte, con $\tau_H = t_H - t_0 = 1$ año. Los instrumentos
se dividen en activos $\mathcal{A}$ y pasivos $\mathcal{L}$, con $\mathcal{K} = \mathcal{A} \cup \mathcal{L}$.
Para cada $k \in \mathcal{K}$ se conocen su moneda $\text{mon}(k) \in \{\text{PEN}, \text{USD}\}$, su tipo
$\text{tipo}(k) \in \{\text{caja}, \text{renta variable}, \text{bono fijo}, \text{bono flotante}\}$, su
vencimiento $T_k$, el factor de referencia de sus cupones si es flotante y su valor de mercado inicial
$z_k^0 > 0$. Para escribir las restricciones de ambos lados con una sola expresión se usa
\[
  z_k = \begin{cases} x_k & k \in \mathcal{A},\\ y_k & k \in \mathcal{L}, \end{cases}
  \qquad
  \epsilon_k = \begin{cases} +1 & k \in \mathcal{A},\\ -1 & k \in \mathcal{L}, \end{cases}
\]
de modo que $\Delta\PN_s = \sum_{k} \epsilon_k\, z_k\, r_{k,s}$.

El mercado en la fecha $t$ es un vector $F_t \in \R^{24}$: las curvas cero de cada moneda en los nodos ON,
1, 3, 5, 10 y 20 años, sus cinco forwards entre nodos consecutivos, el tipo de cambio
$S_t$ (soles por dólar) y el índice de renta variable $E_t$. Los escenarios forman el conjunto finito
$\mathcal{S} = \mathcal{S}^B \cup \mathcal{S}^E$, con $|\mathcal{S}^B| = N_B$ escenarios de remuestreo y
$|\mathcal{S}^E|$ escenarios de estrés, y una medida de probabilidad $p_s \ge 0$, $\sum_{s} p_s = 1$.

\begin{center}
\small
\begin{tabular}{lp{0.55\textwidth}l}
\toprule
Símbolo & Parámetro & Valor del caso \\
\midrule
$B_A,\ B_L$ & presupuestos de activos y pasivos & 1,000 y 800 \\
$\underline{\PN}$ & patrimonio neto mínimo & 120 \\
$\overline{\rho}$ & cociente máximo pasivos/activos & 0.85 \\
$\underline{C}$ & caja mínima & 80 \\
$\underline{w}_g,\ \overline{w}_g$ & pesos mínimo y máximo del grupo $g$ & Cuadro~4 del enunciado \\
$c_k$ & costo de reestructuración por sol movido & según instrumento \\
$\alpha$ & nivel de confianza del CVaR & 0.95 \\
$\lambda$ & aversión al riesgo & 0.25 \\
$N_B,\ b,\ m$ & escenarios de remuestreo, meses por bloque, bloques por escenario & 500, 6, 2 \\
$P^B,\ P^E$ & probabilidad conjunta del remuestreo y del estrés & 0.95, 0.05 \\
\bottomrule
\end{tabular}
\end{center}

\subsection{Problema de optimización}
Dada la matriz de resultados $r \in \R^{|\mathcal{K}| \times |\mathcal{S}|}$, el problema es el programa
lineal
\begin{align}
  \max_{x,\,y,\,\Delta^{+},\,\Delta^{-},\,\zeta,\,u}\quad
    & \sum_{s \in \mathcal{S}} p_s \sum_{k \in \mathcal{K}} \epsilon_k\, z_k\, r_{k,s}
      \;-\; \lambda\Big(\zeta + \frac{1}{1-\alpha}\sum_{s \in \mathcal{S}} p_s\, u_s\Big)
      \;-\; \sum_{k \in \mathcal{K}} c_k\,(\Delta_k^+ + \Delta_k^-) \label{eq:lp-obj}\\
  \text{s.a.}\quad
    & u_s \ge -\textstyle\sum_{k} \epsilon_k\, z_k\, r_{k,s} - \zeta\qquad \forall s \in \mathcal{S} \label{eq:lp-cola}\\
    & \textstyle\sum_{k \in \mathcal{A}} x_k = B_A,\qquad \sum_{k \in \mathcal{L}} y_k = B_L \label{eq:lp-presup}\\
    & \underline{w}_g\, B_{l(g)} \le \textstyle\sum_{k \in \mathcal{K}_g} z_k \le \overline{w}_g\, B_{l(g)}
\qquad \forall g \in \mathcal{G} \label{eq:lp-pesos}\\
    & \textstyle\sum_{k \in \mathcal{A}} x_k - \sum_{k \in \mathcal{L}} y_k \ge \underline{\PN},\qquad
      \sum_{k \in \mathcal{L}} y_k \le \overline{\rho}\,\sum_{k \in \mathcal{A}} x_k \label{eq:lp-pn}\\
    & \textstyle\sum_{k \in \mathcal{A}:\ \text{tipo}(k) = \text{caja}} x_k \ge \underline{C} \label{eq:lp-caja}\\
    & z_k - z_k^0 = \Delta_k^+ - \Delta_k^-\qquad \forall k \in \mathcal{K} \label{eq:lp-reest}\\
    & z_k,\ \Delta_k^+,\ \Delta_k^- \ge 0\quad \forall k,\qquad u_s \ge 0\quad \forall s,\qquad \zeta \in \R.
\end{align}
Cada restricción de pesos $g \in \mathcal{G}$ pertenece a un lado $l(g) \in \{A, L\}$ y actúa sobre un
conjunto de miembros $\mathcal{K}_g$ de ese lado:
\[
\begin{aligned}
  \mathcal{K}_g^{\text{mon} = c} &= \{k : \text{mon}(k) = c\}, &
  \mathcal{K}_g^{\text{flotante}} &= \{k : k \text{ tiene factor de referencia}\},\\
  \mathcal{K}_g^{\text{renta fija}} &= \{k : \text{tipo}(k) \in \{\text{bono fijo}, \text{bono flotante}\}\}, &
  \mathcal{K}_g^{\le 12} &= \{k : M_k(t_0) \le 12\},\\
  & & \mathcal{K}_g^{> 36} &= \{k : M_k(t_0) > 36\},
\end{aligned}
\]
donde $M_k(t)$ es el plazo remanente en meses calendario de $k$ en la fecha $t$. Los grupos de caja y renta
variable se definen igual por tipo, y el de concentración genera una restricción $z_k \le \overline{w}_g\,
B_{l(g)}$ por instrumento. Una vez fijados los presupuestos \eqref{eq:lp-presup}, los pesos
$z_k / B_{l}$ tienen denominador constante y \eqref{eq:lp-pesos} es lineal.

\paragraph{Propiedades.} Tres resultados justifican esta forma del problema.
\begin{enumerate}
  \item \emph{CVaR lineal.} Para $L_s = -\Delta\PN_s$, \textcite{rockafellar2000} muestran que
  \[
    \CVaR_\alpha(L) = \min_{\zeta \in \R}\ \Big\{\zeta + \frac{1}{1-\alpha}\sum_s p_s\,(L_s - \zeta)^+\Big\},
  \]
  y que el conjunto de minimizadores contiene al $\VaR_\alpha(L)$. Las variables $u_s$ con
  \eqref{eq:lp-cola} y $u_s \ge 0$ representan $(L_s - \zeta)^+$ de forma exacta en el óptimo, porque el
  objetivo penaliza $u_s$ con coeficiente $\lambda p_s/(1-\alpha) \ge 0$. Para $\lambda > 0$, maximizar
  \eqref{eq:lp-obj} equivale a maximizar $\E[\Delta\PN] - \lambda\,\CVaR_\alpha(L) - TC$ sobre las carteras
  factibles.
  \item \emph{Valor absoluto sin binarias.} Si $c_k > 0$, en todo óptimo $\min(\Delta_k^+, \Delta_k^-) = 0$:
  si ambas fueran positivas, restar $\min(\Delta_k^+, \Delta_k^-)$ a cada una mantiene \eqref{eq:lp-reest} y
  mejora el objetivo en $2 c_k \min(\Delta_k^+, \Delta_k^-)$. Por lo tanto
  $\Delta_k^+ + \Delta_k^- = |z_k - z_k^0|$ y $TC = \sum_k c_k\, |z_k - z_k^0|$.
  \item \emph{Linealidad en las posiciones.} El resultado por sol invertido $r_{k,s}$ no depende de $z_k$
  (supuesto de escala, Anexo~\ref{app:supuestos}), así que $\Delta\PN_s$ es lineal en $(x, y)$. El problema tiene
  $3|\mathcal{K}| + |\mathcal{S}| + 1$ variables y su conjunto factible es un poliedro. Si es no vacío, el
  objetivo está acotado porque las posiciones lo están por los presupuestos.
\end{enumerate}
Las restricciones \eqref{eq:lp-pn} son redundantes con los presupuestos del caso
($\PN_0 = B_A - B_L = 200$ y $B_L/B_A = 0.80$), pero se mantienen para que el modelo siga siendo válido si
los presupuestos cambian. Con $\lambda = 0$ el problema maximiza el resultado esperado neto de costos y su
solución es, en general, un vértice del poliedro; al crecer $\lambda$ la solución se acerca a la de mínimo
$\CVaR_\alpha$.

\subsection{Modelo de escenarios}\label{app:form-escenarios}
\paragraph{Cambios mensuales.} Para los meses $t = 1, \dots, T$ de la historia anterior a $t_0$, el cambio
del factor $f$ es
\[
  \Delta F_{t,f} = \begin{cases}
    F_{t,f} - F_{t-1,f} & f \text{ tasa (cambio absoluto)},\\
    \ln F_{t,f} - \ln F_{t-1,f} & f \in \{S, E\}\ \text{(log-retorno)}.
  \end{cases}
\]

\paragraph{Remuestreo por bloques.} Para cada inicio $v = 1, \dots, T-b+1$ se define la suma de una
ventana de $b$ meses consecutivos y su versión centrada,
\[
  W_v = \sum_{t=v}^{v+b-1} \Delta F_t, \qquad
  \tilde W_v = W_v - \frac{1}{T-b+1}\sum_{v'=1}^{T-b+1} W_{v'}.
\]
El escenario $s \in \mathcal{S}^B$ suma $m$ ventanas con inicios independientes y uniformes y una deriva
anual $\mu$:
\begin{equation}
  \Delta F_s = \sum_{q=1}^{m} \tilde W_{v_{s,q}} + \mu, \qquad
  v_{s,q} \overset{\text{iid}}{\sim} \mathcal{U}\{1, \dots, T-b+1\}, \qquad b\,m = 12 .
\end{equation}
Como cada $\tilde W_v$ tiene media cero bajo la distribución uniforme, $\E[\Delta F_s] = \mu$ exactamente.
La deriva es nula en las tasas; en el tipo de cambio es la paridad cubierta de tasas,
$\mu_S = \ln\frac{1 + z^{PEN}_{t_0}(\tau_H)}{1 + z^{USD}_{t_0}(\tau_H)}$; y en la renta variable es la
media histórica anualizada, $\mu_E = \frac{12}{T}\sum_t \Delta F_{t,E}$. Como
$\E[e^{X}] \neq e^{\E[X]}$, después de sumar $\mu_S$ todos los log-shocks cambiarios del remuestreo se
desplazan por una constante,
\begin{equation}
  \Delta \ln S_s \leftarrow \Delta \ln S_s + c, \qquad
  c = \ln\frac{F^{S}}{S_{t_0}} - \ln\Big(\frac{1}{N_B}\sum_{s' \in \mathcal{S}^B} e^{\Delta \ln S_{s'}}\Big),
\end{equation}
donde $F^{S}/S_{t_0} = \frac{1 + z^{PEN}_{t_0}(\tau_H)}{1 + z^{USD}_{t_0}(\tau_H)}$. Así la media de
$S_{t_H}/S_{t_0}$ en la muestra es exactamente la forward de paridad. Cada escenario de remuestreo tiene
probabilidad $p_s = P^B / N_B$.

\paragraph{Escenarios de estrés.} Cada escenario $s \in \mathcal{S}^E$ define, por moneda, un
desplazamiento paralelo $\delta_s$ y, opcionalmente, desplazamientos $\delta_s^{\text{corto}}$ y
$\delta_s^{\text{largo}}$ que reemplazan al paralelo en los nodos de cada tramo,
$\mathcal{N}^{\text{corto}} = \{\text{ON}, 1\}$ y $\mathcal{N}^{\text{largo}} = \{10, 20\}$. En un nodo de
transición $t$ entre los nodos fijos vecinos $t_a < t < t_b$, con desplazamientos $\delta_a$ y $\delta_b$, se
interpola linealmente el producto $t\,\delta(t)$:
\begin{equation}
  t\,\delta(t) = t_a \delta_a + (t_b \delta_b - t_a \delta_a)\,\frac{t - t_a}{t_b - t_a}.
\end{equation}
Con capitalización continua, $t\,z(t)$ es el logaritmo del factor de capitalización, de modo que esta regla
reparte por igual entre los segmentos de la transición el cambio de la forward media
$\frac{t_b \delta_b - t_a \delta_a}{t_b - t_a}$; con capitalización efectiva anual la igualdad es aproximada.
Ese cambio medio lo fijan los datos del escenario y ninguna regla de transición puede evitarlo. En tipo de
cambio y renta variable, $\Delta \ln S_s = \ln(1 + \phi_s)$ y $\Delta \ln E_s = \ln(1 + \eta_s)$, con
$\phi_s$ y $\eta_s$ las variaciones porcentuales del escenario, sin deriva. Cada estrés recibe su
probabilidad, y $\sum_{s \in \mathcal{S}^E} p_s = P^E$.

\paragraph{Coherencia del mercado.} En todo escenario, las forwards se recalculan como implícitas de la
curva cero desplazada,
\[
  f(t_1, t_2) = \left(\frac{(1 + z(t_2))^{t_2}}{(1 + z(t_1))^{t_1}}\right)^{1/(t_2 - t_1)} - 1,
\]
con $z(\cdot)$ interpolada linealmente entre nodos. No se impone un piso a las tasas.

\subsection{Valorización y resultado por instrumento}\label{app:form-revalorizacion}
\paragraph{Descuento y calibración.} Un instrumento con flujos $\{CF_{k,n}\}$ en las fechas $\tau_n$ y moneda
$c$ vale, en la fecha $t$ y con el mercado $F$,
\[
  V_k(t; F) = S_c(F) \sum_{n:\ \tau_n > t} CF_{k,n}\,\big(1 + z^{c}_F(\tau_n - t) + \sigma_k\big)^{-(\tau_n - t)},
\]
donde $S_c$ convierte a soles ($S_{\text{PEN}} = 1$). El spread $\sigma_k$ es la única solución de
$V_k(t_0; F_{t_0}) = z_k^0$ (el valor presente es estrictamente decreciente en $\sigma_k$) y se mantiene
constante hasta el horizonte.

\paragraph{Trayectoria dentro del año.} Entre $t_0$ y $t_H$ el mercado del escenario $s$ se mueve
linealmente:
\[
  F_s(\tau) = F_{t_0} + \theta(\tau)\,\Delta F_s, \qquad \theta(\tau) = \frac{\tau - t_0}{t_H - t_0} \in [0, 1].
\]
Un cupón flotante con fecha de fijación $\tau^{\text{fix}}$ y spread contractual $\kappa_k$ paga, sobre el
nocional $N_k$ y con frecuencia $\omega_k$,
\[
  CF_{k,n} = \frac{N_k}{\omega_k}\,\big(R_k(\tau^{\text{fix}}_n) + \kappa_k\big), \qquad
  R_k(\tau) = \begin{cases}
    R_k^{\text{base}}(\tau) & \tau \le t_0,\\
    R_k\big(F_s(\min(\tau, t_H))\big) & \tau > t_0,
  \end{cases}
\]
donde $R_k^{\text{base}}$ es la proyección base, que no cambia para los cupones ya fijados, y $R_k(F)$ es el
valor del factor de referencia de $k$ en el mercado $F$.

\paragraph{Resultado al horizonte.} El resultado por sol invertido es
\begin{equation}
  r_{k,s} = \underbrace{\frac{V_k(t_H; F_s(t_H))}{z_k^0} - 1}_{\text{revalorización}}
  \;+\; \underbrace{\frac{1}{z_k^0} \sum_{n:\ t_0 < \tau_n \le t_H} S_c\big(F_s(\tau_n)\big)\, CF_{k,n}}_{\text{flujos del año}},
\end{equation}
donde los flujos del año se convierten al tipo de cambio de su fecha de pago y no se reinvierten. La caja y
la renta variable no tienen flujos:
\[
  r_{k,s}^{\text{caja}} = \Big(1 + \big(z^c_{t_0}(\text{ON}) + \tfrac{1}{2}\Delta z^c_s(\text{ON})\big)\,\tau_H\Big)\frac{S_c(F_s(t_H))}{S_c(F_{t_0})} - 1,
  \qquad
  r_{k,s}^{\text{RV}} = e^{\Delta \ln E_s}\,\frac{S_c(F_s(t_H))}{S_c(F_{t_0})} - 1 .
\]
La caja devenga la tasa ON media de la trayectoria lineal en tasa simple, y la renta variable no paga
dividendos. Para un pasivo, un $r_{k,s}$ mayor (más intereses pagados o un mayor valor de la deuda) reduce
$\Delta\PN_s$ por el signo $\epsilon_k = -1$.

\subsection{Medidas que se reportan}
\paragraph{VaR y CVaR \emph{ex post}.} Con $L_s = -\Delta\PN_s$ y la distribución discreta $\{(L_s, p_s)\}$,
\begin{align}
  \VaR_\alpha &= \min\Big\{\ell : \sum_{s:\ L_s \le \ell} p_s \ge \alpha\Big\},\\
  \CVaR_\alpha &= \frac{1}{1-\alpha}\Big[\sum_{L_s > \VaR_\alpha} p_s L_s
    + \VaR_\alpha\Big(1-\alpha-\sum_{L_s > \VaR_\alpha} p_s\Big)\Big].
\end{align}
El segundo término asigna al átomo del $\VaR$ solo la fracción de su probabilidad que completa la cola de
masa $1-\alpha$. Se calcula así, y no leyendo $\zeta$, porque con $\lambda = 0$ el valor de $\zeta$ es
arbitrario y con átomos de probabilidad el cuantil puede no ser único.

\paragraph{Resultados esperados.}
\[
  \E[\Delta\PN] = \sum_s p_s\,\Delta\PN_s,\qquad
  \E[\PN_H] = \PN_0 - TC + \E[\Delta\PN],
\]
\[
  \bar r_A = \frac{\sum_{k \in \mathcal{A}} x_k\, \bar r_k}{B_A},\qquad
  \bar r_L = \frac{\sum_{k \in \mathcal{L}} y_k\, \bar r_k}{B_L},
\]
con $\bar r_k = \sum_s p_s r_{k,s}$. La volatilidad es la desviación estándar de $\Delta\PN_s$ bajo $p$.

\paragraph{Restricciones al horizonte.} Las restricciones se imponen en $t_0$, sobre la estructura que sale
de la decisión, y en $t_H$ se reporta la probabilidad de incumplirlas. La caja al horizonte acumula los
flujos netos del año y descuenta el costo de transacción pagado en $t_0$:
\[
  C_{H,s} = \sum_{k:\ \text{tipo}(k) = \text{caja}} x_k\,(1 + r_{k,s})
  + \sum_{k} \epsilon_k\, \frac{z_k}{z_k^0} \sum_{t_0 < \tau_n \le t_H} S_c(F_s(\tau_n))\, CF_{k,n} - TC,
\]
lo que asegura que $A_{H,s} - L_{H,s} = \PN_0 - TC + \Delta\PN_s$. Para cada escenario se verifican
$\PN_{H,s} \ge \underline{\PN}$, $L_{H,s} \le \overline{\rho}\,A_{H,s}$, $C_{H,s} \ge \underline{C}$ y cada
grupo \eqref{eq:lp-pesos} con las tenencias en $t_H$ y la pertenencia $M_k(t_H)$. Un incumplimiento es
\emph{estructural} si ocurre aun con los montos de $t_0$, es decir, por el solo paso del tiempo, y \emph{de
mercado} si se debe a la variación de los valores.

\paragraph{Variantes de comparación.} La cartera de mínima varianza resuelve, con las mismas restricciones y
escenarios, el programa cuadrático
\[
  \min_{x,\,y,\,\Delta^\pm}\ \sum_s p_s\big(\Delta\PN_s - \E[\Delta\PN]\big)^2
  \quad\text{s.a.}\quad \E[\Delta\PN] - TC \ge \E[\Delta\PN^*] - TC^*,
\]
donde el lado derecho corresponde a la solución de \eqref{eq:lp-obj}. La variante de vencimientos al horizonte
agrega a \eqref{eq:lp-pesos} los grupos de vencimiento con la pertenencia $M_k(t_H)$, y su costo es la
diferencia de valor óptimo frente al problema base.

\subsection{Supuestos y decisiones de modelación}\label{app:supuestos}
\paragraph{Alcance de la decisión.}
\begin{enumerate}
  \item \emph{Decisión estática.} Se decide una vez en $t_0$ y no se rebalancea dentro del año. Es lo que pide el
  caso y mantiene el problema en una sola etapa.
  \item \emph{Solo los instrumentos existentes.} Se optimizan los seis activos y los seis pasivos; la deuda nueva
  no entra, porque el caso asignado se limita a reestructurar el balance existente.
  \item \emph{Escala.} Cambiar $z_k$ equivale a comprar o emitir más del mismo contrato, con el mismo spread y sin
  impacto de precio. Este supuesto hace que $r_{k,s}$ no dependa de $z_k$ y el problema sea lineal.
  \item \emph{Sin posiciones cortas}: $z_k \ge 0$, porque el caso no ofrece esa alternativa.
\end{enumerate}

\paragraph{Objetivo y costos.}
\begin{enumerate}\setcounter{enumi}{4}
  \item \emph{Riesgo como penalización.} El CVaR entra en el objetivo ponderado por $\lambda$, como pide el
  enunciado. Imponerlo como restricción ($\CVaR \le c$) y variar $c$ traza la misma frontera eficiente que
  variar $\lambda$; maximizar el cociente $\E/\CVaR$, en cambio, entrega un solo punto de ella.
  \item \emph{Costo simétrico y único.} Solo se aplica el costo de reestructuración, a subidas y bajadas por igual.
  Sumarle el costo de prepago a las reducciones de pasivos cobraría dos veces el mismo movimiento.
  \item \emph{Costos fuera de la pérdida.} $TC$ es cierto y se paga en $t_0$: se resta una sola vez en el objetivo y
  no entra en $L_s$, de modo que el CVaR mide solo el riesgo de mercado.
  \item \emph{Límites en $t_0$.} Las restricciones se imponen sobre la estructura que sale de la decisión, como
  dice el enunciado, y en $t_H$ se reportan. Imponerlas también al horizonte obligaría a cubrir la deriva de
  mercado de todos los límites; en su lugar se cuantifica el costo de exigir los vencimientos en $t_H$.
\end{enumerate}

\paragraph{Escenarios.}
\begin{enumerate}\setcounter{enumi}{8}
  \item \emph{Bloques de seis meses.} Los cambios mensuales son muy persistentes: la autocorrelación de orden uno
  es 0.94 en la tasa ON en soles, 0.96 en la de dólares, 0.91 en la renta variable y 0.56 en el tipo de cambio.
  Doce meses independientes subestimarían la varianza anual. Dos bloques de seis meses conservan esa persistencia
  y ofrecen $(T-b+1)^2$ combinaciones, frente a solo $T-11$ con un bloque único de doce meses.
  \item \emph{Centrar las ventanas, no los meses.} Sin vuelta circular, los meses de los extremos aparecen en menos
  ventanas, y centrar los meses dejaría a $\Delta F_s$ con media distinta de $\mu$. Unir el fin de la muestra con
  su inicio (remuestreo circular) crearía saltos artificiales.
  \item \emph{Deriva mixta.} Centrar todo eliminaría la prima de la renta variable y la del tipo de cambio, y la renta
  variable quedaría en su mínimo permitido. Mantener la media histórica de las tasas proyectaría el ciclo alcista
  de 2021--2025. Por eso las tasas se centran, el tipo de cambio sigue la paridad de tasas y la renta variable
  conserva su media histórica.
  \item \emph{Ajuste del tipo de cambio a la forward.} Sin el desplazamiento $c$, la convexidad de $e^{X}$ y el error
  de muestreo de la media dejarían el tipo de cambio esperado por encima de la forward, a favor de estar largo en
  dólares.
  \item \emph{Tramos del estrés.} El enunciado no define los nodos del tramo corto ni del largo; se toman
  $\{\text{ON}, 1\}$ y $\{10, 20\}$. Los desplazamientos de cada tramo se interpretan como cambios que reemplazan al
  paralelo, no como niveles de tasa: leídos como nivel, un 4\,\% casi no movería el tramo corto en soles, lo que
  contradice la descripción del escenario.
  \item \emph{Transición por $t\,\delta(t)$.} Un escalón entre tramos, o promediar los desplazamientos en los nodos
  intermedios, produce forwards que se mueven en sentido contrario a las tasas. Interpolar $t\,\delta(t)$ reparte el
  cambio de forward que imponen los datos de manera pareja.
  \item \emph{Forwards implícitas y sin piso.} Las forwards se derivan de la curva cero desplazada para que el
  mercado de cada escenario no admita arbitraje entre plazos. No se impone un piso a las tasas, para no
  truncar la distribución; la tasa mínima de cada escenario se revisa como control.
\end{enumerate}

\paragraph{Valorización.}
\begin{enumerate}\setcounter{enumi}{15}
  \item \emph{Revalorización completa.} Los resultados se obtienen revalorizando los flujos de cada instrumento, no
  con duración o DV01, para capturar la convexidad y los efectos no paralelos de los escenarios.
  \item \emph{Trayectoria lineal.} Una misma regla para la caja, los cupones flotantes y la conversión de flujos en
  dólares evita tratar los activos y los pasivos de forma asimétrica dentro del año.
  \item \emph{Cupón flotante al nivel del factor de referencia.} Cada cupón paga el valor del factor que indexa el
  contrato en su fecha de fijación, más el spread contractual. Es la convención con la que se construyó la
  proyección base del caso. La alternativa, la forward implícita del período de cada cupón, se evalúa como
  sensibilidad (Sección~\ref{sec:sensibilidad}).
  \item \emph{Spread constante.} No hay datos sobre la dinámica de los spreads de crédito; el spread se calibra
  al valor de mercado en $t_0$ y se mantiene fijo.
  \item \emph{Sin reinversión ni dividendos.} Los flujos del año quedan en caja sin rendir intereses, y la renta
  variable no paga dividendos. Ambos efectos son de segundo orden frente a los shocks del escenario.
\end{enumerate}

\paragraph{Alternativas descartadas.} Remuestreo independiente mes a mes (subestima la varianza anual);
remuestreo estacionario con bloques de largo aleatorio (más difícil de verificar); simulación paramétrica
desde la covarianza, que es singular porque las forwards se derivan de las curvas cero; el estrés como
restricción dura; costos fijos de transacción, que requerirían variables binarias; y el CVaR a varios
niveles de confianza a la vez.
